\begin{itemize}
\item \textbf{MIPLIB 2017}~\citep{gleixner2021miplib}: 80 of the 240 benchmark-set instances, chosen
deterministically for diversity --- at most one instance per MIPLIB group, at most 2.5M non-zeros, not
infeasible, stratified by MIPLIB status and size tercile. Budget $T=60$\,s. Reference: published optimum or best
known value. The set is 96\,\% ``easy'' by MIPLIB's definition (solved within an hour by some solver), which is
not the same as easy within 60\,s: only 56 of 80 reach a 1\,\% gap within 60\,s under any strategy.
\item \textbf{ML4CO item placement} and \textbf{load balancing}~\citep{gasse2022ml4co}: two homogeneous
generated families from the NeurIPS 2021 competition. Training instances 0--74 of each family's official training
split (cuOpt) / 0--99 (CPU strategies), and the first 25 instances of the official validation split as test
instances. Budget $T=30$\,s. Reference: best value found by any of our runs.
\item \textbf{PGLib-UC}~\citep{pglibuc}: all 56 unit-commitment cases (three base systems: CAISO, FERC, RTS-GMLC),
built with the benchmark's own reference Pyomo model~\citep{knueven2020uc} (about 250k rows and 250k columns,
40\,\% binary). Budget $T=60$\,s. Split by base system.
\end{itemize}
The budgets are shorter than the spec's 300\,s screening: the machine has one GPU, cuOpt is the throughput
bottleneck, and one night was available. The 10\,s and 30\,s views are cut from the same trajectories. On MIPLIB,
HiGHS and SCIP defaults were additionally run for 300\,s to place instances into the spec's
SHORT/MEDIUM/HARD bands (Appendix).
